\documentclass[a4paper,12pt]{ltjsarticle}
\usepackage[margin=25mm]{geometry}
\usepackage{amsmath,amssymb}
\usepackage{enumitem}

\begin{document}

\begin{center}
{\LARGE 数I・A 共通テスト風 解答（5-a）}
\end{center}

\vspace{1em}

\section*{第1問（数と式・二次関数）}
\[
f(x)=x^2-6x+a=(x-3)^2+(a-9)
\]
\begin{enumerate}[label=\arabic*.]
  \item 最小値は \(\,a-9\)（\(x=3\) のとき）。
  \item 判別式より
  \[
  D=36-4a>0 \iff a<9.
  \]
  \item \(\alpha<1<\beta\) は \(f(1)<0\) と同値。
  \[
  f(1)=1-6+a=a-5<0 \iff a<5.
  \]
  また \(\alpha,\beta\) が実数であるには \(a\le 9\) が必要だが、\(a<5\) なら自動的に満たす。  
  よって \(\,a<5\)。
  \item \(a=5\) のとき
  \[
  x^2-6x+5\le0
  \iff (x-1)(x-5)\le0
  \iff 1\le x\le5.
  \]
\end{enumerate}

\section*{第2問（場合の数・確率）}
全事象は
\[
{8 \choose 3}=56
\]
通り。
\begin{enumerate}[label=\arabic*.]
  \item 偶数4枚・奇数4枚より、和が偶数になるのは  
  （偶偶偶）または（偶奇奇）。
  \[
  {4 \choose 3}+{4 \choose 1}{4 \choose 2}=4+24=28
  \]
  よって確率は
  \[
  \frac{28}{56}=\frac12.
  \]
  \item 最大が6とは、\(\{1,2,3,4,5\}\)から2枚選び6を加える場合：
  \[
  {5 \choose 2}=10
  \]
  よって
  \[
  \frac{10}{56}=\frac{5}{28}.
  \]
  \item 3の倍数を含まない場合を補集合で数える。  
  3の倍数は \(\{3,6\}\) の2枚、非3の倍数は6枚。
  \[
  P(\text{積が3の倍数})=1-\frac{{6 \choose 3}}{{8 \choose 3}}
  =1-\frac{20}{56}
  =\frac{36}{56}
  =\frac{9}{14}.
  \]
  \item 条件 \(a+b>c\) を満たす組（\(1\le a<b<c\le 8\)）を数えると18通り。  
  よって
  \[
  \frac{18}{56}=\frac{9}{28}.
  \]
\end{enumerate}

\section*{第3問（図形と計量）}
\[
AB=6,\ AC=8,\ \angle A=60^\circ
\]
\begin{enumerate}[label=\arabic*.]
  \item 面積
  \[
  S=\frac12\cdot6\cdot8\cdot\sin60^\circ
  =24\cdot\frac{\sqrt3}{2}
  =12\sqrt3.
  \]
  \item 余弦定理で
  \[
  BC^2=6^2+8^2-2\cdot6\cdot8\cos60^\circ=52
  \]
  より \(BC=2\sqrt{13}\)。  
  正弦定理 \( \dfrac{a}{\sin A}=2R \)（\(a=BC\)）から
  \[
  R=\frac{a}{2\sin A}
  =\frac{2\sqrt{13}}{2\cdot(\sqrt3/2)}
  =\frac{2\sqrt{13}}{\sqrt3}
  =\frac{2\sqrt{39}}{3}.
  \]
  \item \(BC\) を底辺、高さを \(h\) とすると
  \[
  12\sqrt3=\frac12\cdot(2\sqrt{13})\cdot h
  \]
  よって
  \[
  h=\frac{12\sqrt3}{\sqrt{13}}
  =\frac{12\sqrt{39}}{13}.
  \]
\end{enumerate}

\section*{第4問（データの分析）}
\[
H=50+10\frac{x-64}{12}
\]
\begin{enumerate}[label=\arabic*.]
  \item \(x=76\) のとき
  \[
  H=50+10\cdot\frac{12}{12}=60.
  \]
  \item
  \[
  50+10\frac{x-64}{12}\ge60
  \iff \frac{x-64}{12}\ge1
  \iff x\ge76.
  \]
\end{enumerate}

\end{document}
